\documentclass[12pt,a4paper]{article}

% Packages
\usepackage[utf8]{inputenc}
\usepackage[T1]{fontenc}
\usepackage{amsmath,amssymb}
\usepackage{geometry}
\geometry{margin=1in}
\usepackage{setspace}
\onehalfspacing
\usepackage{natbib}
\usepackage{hyperref}
\hypersetup{colorlinks=true,linkcolor=blue,citecolor=blue,urlcolor=blue}
\usepackage{booktabs}
\usepackage{graphicx}
\usepackage{enumitem}

% Title
\title{\textbf{Corporate AI Accessibility and Information Asymmetry: \\
Evidence from Firms' Robots.txt Policies}\\
\vspace{0.5cm}
\large Research Proposal}

\author{}
\date{\today}

\begin{document}
\maketitle

% ============================================================
\section{Introduction and Motivation}
% ============================================================

The rise of large language models (LLMs) and AI-powered information processing tools has fundamentally altered how investors, analysts, and regulators acquire and analyze corporate information. As of 2025, an estimated 85\% of downloads from the SEC's EDGAR database originate from automated bots rather than human readers.\footnote{See the SEC speech by Scott Bauguess, ``The Role of Machine Readability in an AI World,'' May 3, 2018, available at \url{https://www.sec.gov/newsroom/speeches-statements/speech-bauguess-050318}. The 85\% figure refers to EDGAR traffic data discussed therein.} Financial analysts increasingly employ AI tools to process earnings calls, 10-K filings, and corporate press releases, while retail investors interact with firm information through AI-powered search engines and chatbots.

A growing body of literature examines how firms respond to AI readership on the \emph{supply side} of disclosure. \citet{CaoJiangYangZhang2023}, published in the \emph{Review of Financial Studies} (Vol.\ 36, No.\ 9, pp.\ 3603--3642), document that firms actively adjust their filings to be more machine-friendly and avoid words that NLP algorithms perceive as negative. On a related front, \citet{BlankespoorDeHaanLi2024}, in a working paper available on SSRN (Abstract ID: 4986017, last revised 2024), find that up to 4.5\% of new disclosure text across earnings press releases, conference call remarks, risk factors, MD\&As, and S-1 filings was generated by AI tools by 2024.

However, an equally important but entirely unstudied dimension exists on the \emph{demand side}: \textbf{whether firms allow AI systems to access their information in the first place.} Every website operator can configure a \texttt{robots.txt} file---a publicly observable, machine-readable protocol---that specifies which automated crawlers are permitted or denied access. Since OpenAI introduced its GPTBot web crawler in August 2023, firms face an explicit, binary choice: allow or block AI systems from reading their web content.\footnote{OpenAI announced GPTBot and its \texttt{robots.txt} compliance on August 8, 2023. See \url{https://platform.openai.com/docs/bots}.} Industry data show that the share of top-1000 websites blocking GPTBot rose from approximately 5\% at introduction to 35.7\% by mid-2025.\footnote{Paul Calvano, ``AI Bots and Robots.txt,'' August 21, 2025, available at \url{https://paulcalvano.com/2025-08-21-ai-bots-and-robots-txt/}. These figures pertain to the top 1,000 global websites and may not directly generalize to publicly listed firms.}

Crucially, \textbf{no existing study examines publicly listed firms' AI-blocking behavior or its financial consequences.} This proposal outlines a research agenda to fill this gap by constructing a novel measure of \emph{corporate AI accessibility} and investigating its implications for information asymmetry, analyst coverage, and capital market outcomes.

% ============================================================
\section{Research Questions}
% ============================================================

This project addresses three primary research questions:

\begin{enumerate}[label=\textbf{RQ\arabic*.}]
    \item \textbf{Who blocks AI?} What firm characteristics predict the decision to block AI crawlers (e.g., GPTBot, ClaudeBot, CCBot) via \texttt{robots.txt}? Do firms with greater information asymmetry, weaker governance, or higher proprietary information costs exhibit systematically different AI-blocking behavior?

    \item \textbf{What are the capital market consequences?} Does blocking AI crawlers affect bid-ask spreads, analyst forecast dispersion, trading volume, or post-earnings-announcement drift (PEAD)?

    \item \textbf{Does AI accessibility affect firm valuation?} Drawing on \citet{Merton1987}'s investor recognition hypothesis,\footnote{Robert C.\ Merton, ``A Simple Model of Capital Market Equilibrium with Incomplete Information,'' \emph{Journal of Finance}, Vol.\ 42, No.\ 3, July 1987, pp.\ 483--510.} does restricting AI access reduce investor recognition in an era where AI mediates information acquisition?
\end{enumerate}

% ============================================================
\section{Relation to Existing Literature}
% ============================================================

This project contributes to and bridges several literatures. Table~\ref{tab:positioning} summarizes the positioning relative to the most directly related papers.

\begin{table}[htbp]
\centering
\caption{Positioning Relative to Existing Literature}
\label{tab:positioning}
\small
\begin{tabular}{p{4.5cm}p{3.5cm}p{5cm}}
\toprule
\textbf{Paper} & \textbf{Focus} & \textbf{Distinction from This Project} \\
\midrule
\citet{CaoJiangYangZhang2023}, \emph{RFS} 36(9), pp.\ 3603--3642 & Firms adapt disclosure \emph{language} to AI readers & We study whether firms \emph{block} AI readers entirely \\
\addlinespace
\citet{BlankespoorDeHaanLi2024}, SSRN WP \#4986017 & Firms use GenAI to \emph{write} disclosures & We study whether firms let AI \emph{read} their information \\
\addlinespace
\citet{ChangHe2025}, SSRN WP \#5159436 & Legal liability of \texttt{robots.txt} violations & We study \emph{financial} consequences of \texttt{robots.txt} choices \\
\addlinespace
\citet{Reitinger2025}, SSRN WP \#5555040 & \texttt{robots.txt} usage across \emph{general} websites & We focus on \emph{publicly listed firms} specifically \\
\addlinespace
SEC/XBRL literature (e.g., FDTA 2022) & Machine readability of \emph{mandatory} filings & We study machine accessibility of \emph{voluntary} information \\
\bottomrule
\end{tabular}
\end{table}

\subsection{Corporate Disclosure and Information Processing}

The literature on disclosure readability has long examined how the complexity of corporate filings affects information asymmetry. \citet{Li2008} (\emph{The Accounting Review}, Vol.\ 83, No.\ 1, pp.\ 221--236) introduced the Fog Index as a measure of annual report readability and found that less readable 10-K filings are associated with lower earnings and greater earnings persistence. Our project extends this concept from \emph{human} readability to \emph{machine} readability---specifically, whether AI systems can access corporate information at all.

\subsection{AI and Corporate Finance}

A rapidly expanding literature investigates AI's role in corporate finance. \citet{CaoYangZhao2024} (``Information in Disclosing Emerging Technologies: Evidence from AI Disclosure,'' SSRN WP \#4987085) examine the information content of voluntary AI disclosures in annual reports. Our work complements this stream by studying firms' revealed preferences toward AI through their \texttt{robots.txt} choices---an observable action rather than a textual disclosure.

\subsection{Information Asymmetry and Investor Recognition}

By potentially raising the cost of AI-mediated information acquisition, blocking AI crawlers may reduce investor recognition in the sense of \citet{Merton1987}. This connects to the broader literature on information frictions and asset pricing, including work on analyst coverage and information intermediation.

% ============================================================
\section{Data and Variable Construction}
% ============================================================

\subsection{AI Accessibility Variable}

The primary independent variable, \textbf{AI Accessibility}, is constructed by scraping the \texttt{robots.txt} file from each sample firm's corporate website. Specifically, for each firm $i$ at time $t$, we define:

\begin{equation}
\text{AIBlock}_{i,t} = \begin{cases} 1 & \text{if firm } i \text{'s robots.txt disallows GPTBot (or other AI crawlers) at time } t \\ 0 & \text{otherwise} \end{cases}
\end{equation}

We also construct a continuous measure:

\begin{equation}
\text{AIAccessIndex}_{i,t} = 1 - \frac{\text{Number of AI crawlers blocked}_{i,t}}{\text{Total number of major AI crawlers}}
\end{equation}

\noindent where the set of major AI crawlers includes GPTBot (OpenAI), ClaudeBot (Anthropic), Google-Extended (Google), CCBot (Common Crawl), and others as they emerge.

Additionally, we construct a measure based on \emph{actual fetch success}:
\begin{equation}
\text{FetchSuccess}_{i,t} = \frac{\text{Number of successfully fetched key pages}_{i,t}}{\text{Total key pages attempted}_{i,t}}
\end{equation}
\noindent where key pages include the investor relations page, most recent press releases, and corporate governance section. This captures not only \texttt{robots.txt} directives but also other technical barriers (JavaScript rendering, CAPTCHAs, rate limiting).

\subsection{Data Sources and Sample}

\begin{itemize}
    \item \textbf{Robots.txt data}: Scraped monthly from corporate websites of S\&P 500 and Russell 3000 firms. Historical snapshots can be retrieved from the Wayback Machine (\url{https://web.archive.org/}) to construct a panel from August 2023 (GPTBot introduction) onward.

    \item \textbf{Firm financials}: Compustat (via WRDS) for standard firm characteristics.

    \item \textbf{Stock market data}: CRSP for returns, volume, and bid-ask spreads.

    \item \textbf{Analyst data}: I/B/E/S for analyst coverage, forecast dispersion, and earnings surprises.

    \item \textbf{Institutional ownership}: Thomson Reuters 13F filings.

    \item \textbf{Corporate governance}: ISS (formerly RiskMetrics) for board characteristics.

    \item \textbf{Corporate website URLs}: Mapped from SEC EDGAR company filings or Compustat's web address field.
\end{itemize}

\subsection{Sample Period and Construction}

The natural starting point is \textbf{August 2023}, when OpenAI introduced GPTBot with \texttt{robots.txt} compliance. The sample extends through the most recent available data. The staggered introduction of additional AI crawlers (Anthropic's ClaudeBot, Google's AI-related crawlers) provides further temporal variation.

% ============================================================
\section{Empirical Strategy}
% ============================================================

\subsection{Cross-Sectional Determinants (RQ1)}

We first examine the determinants of AI-blocking behavior:

\begin{equation}
\text{AIBlock}_{i,t} = \alpha + \beta \mathbf{X}_{i,t-1} + \gamma_j + \delta_t + \varepsilon_{i,t}
\end{equation}

\noindent where $\mathbf{X}_{i,t-1}$ includes lagged firm size, book-to-market ratio, institutional ownership, analyst coverage, R\&D intensity (as a proxy for proprietary information), litigation risk, information asymmetry measures (e.g., bid-ask spread, PIN), and governance characteristics. $\gamma_j$ and $\delta_t$ denote industry and time fixed effects, respectively.

\subsection{Capital Market Consequences (RQ2)}

To estimate the causal effect of AI blocking on information asymmetry, we employ several approaches:

\paragraph{Panel Regressions with Firm Fixed Effects.}
\begin{equation}
Y_{i,t} = \alpha_i + \beta \cdot \text{AIBlock}_{i,t} + \gamma \mathbf{X}_{i,t} + \delta_t + \varepsilon_{i,t}
\end{equation}

\noindent where $Y_{i,t}$ is a measure of information asymmetry (bid-ask spread, Amihud illiquidity, analyst forecast dispersion), $\alpha_i$ are firm fixed effects, and identification comes from within-firm variation in AI-blocking status over time.

\paragraph{Event Study around GPTBot Introduction.}
We exploit the introduction of GPTBot in August 2023 as a natural experiment. Before August 2023, no firm could block GPTBot because it did not exist. The cross-sectional variation in firms' \emph{responses} to this new technology---some block immediately, others do not---provides quasi-experimental variation. We estimate:

\begin{equation}
Y_{i,t} = \alpha_i + \sum_{\tau = -K}^{K} \beta_\tau \cdot \mathbf{1}[t = \text{Aug 2023} + \tau] \times \text{EarlyBlocker}_{i} + \gamma \mathbf{X}_{i,t} + \delta_t + \varepsilon_{i,t}
\end{equation}

\noindent where $\text{EarlyBlocker}_{i}$ equals one if firm $i$ blocked GPTBot within the first three months of its introduction.

\paragraph{Difference-in-Differences around Policy Changes.}
If a firm \emph{switches} from allowing to blocking AI crawlers (or vice versa), we can estimate the within-firm effect using a staggered DID design. Given recent econometric advances on staggered DID bias, we plan to implement the \citet{CallawayS2021} estimator.\footnote{Brantly Callaway and Pedro H.C.\ Sant'Anna, ``Difference-in-Differences with Multiple Time Periods,'' \emph{Journal of Econometrics}, Vol.\ 225, No.\ 2, December 2021, pp.\ 200--230.}

\subsection{Firm Valuation (RQ3)}

We examine whether AI accessibility is priced:
\begin{equation}
\text{Tobin's } q_{i,t} = \alpha + \beta \cdot \text{AIAccessIndex}_{i,t} + \gamma \mathbf{X}_{i,t} + \gamma_j + \delta_t + \varepsilon_{i,t}
\end{equation}

We also test whether the valuation effect is moderated by analyst coverage (following the investor recognition logic).

% ============================================================
\section{Identification Challenges and Mitigation}
% ============================================================

\begin{table}[htbp]
\centering
\caption{Identification Challenges and Proposed Solutions}
\label{tab:identification}
\small
\begin{tabular}{p{4cm}p{8.5cm}}
\toprule
\textbf{Challenge} & \textbf{Mitigation Strategy} \\
\midrule
Endogeneity: opaque firms may both block AI and have higher information asymmetry & (1) Firm fixed effects exploit within-firm variation; (2) Event study around GPTBot introduction; (3) Instrumental variable: industry-level AI-blocking rate excluding firm $i$ \\
\addlinespace
Selection: firm size and IT sophistication correlate with robots.txt usage & Control for firm size, IT investment, website complexity; use matched samples (PSM/CEM) \\
\addlinespace
Confounding: robots.txt may block all bots, not just AI & Distinguish firms that specifically block AI crawlers (GPTBot, ClaudeBot) from those using generic bot-blocking. Many firms allow Googlebot but block GPTBot---this variation is informative \\
\addlinespace
Measurement: robots.txt is not the only barrier & Supplement with actual fetch tests to capture JavaScript rendering, CAPTCHAs, etc. \\
\addlinespace
Reverse causality & Use lagged AI-blocking status; exploit timing of crawler introductions as shocks \\
\bottomrule
\end{tabular}
\end{table}

% ============================================================
\section{Expected Contributions}
% ============================================================

\begin{enumerate}
    \item \textbf{Novel variable.} We introduce a new, directly observable measure of corporate information accessibility in the AI era. Unlike textual readability measures (e.g., Fog Index) that are continuous and subject to measurement error, \texttt{robots.txt} directives are binary, publicly observable, and clearly reflect deliberate firm choices.

    \item \textbf{New dimension of disclosure.} Existing literature focuses on \emph{what} firms disclose and \emph{how} they write it. We study a prior question: whether firms allow AI systems to \emph{access} their information at all. This extends the disclosure literature into the AI era.

    \item \textbf{Policy relevance.} As the SEC continues to expand machine-readable filing requirements (e.g., the Financial Data Transparency Act of 2022), understanding firms' voluntary choices regarding AI accessibility becomes directly policy-relevant. Our findings may inform regulatory discussions about whether AI accessibility of corporate information should be mandated or standardized.

    \item \textbf{Unique, self-constructed dataset.} The dataset does not depend on any commercial data provider and can be constructed entirely from publicly available information, enabling full replication.
\end{enumerate}

% ============================================================
\section{Preliminary Timeline}
% ============================================================

\begin{table}[htbp]
\centering
\begin{tabular}{ll}
\toprule
\textbf{Phase} & \textbf{Timeline} \\
\midrule
Data collection: scrape robots.txt for Russell 3000 & Months 1--2 \\
Historical panel via Wayback Machine (Aug 2023--present) & Months 2--4 \\
Descriptive analysis and variable validation & Months 4--5 \\
Cross-sectional determinants analysis (RQ1) & Months 5--7 \\
Capital market consequences analysis (RQ2) & Months 7--10 \\
Valuation analysis (RQ3) & Months 10--12 \\
Paper writing and conference submissions & Months 12--15 \\
\bottomrule
\end{tabular}
\end{table}

% ============================================================
% References
% ============================================================

\newpage
\begin{thebibliography}{99}

\bibitem[Blankespoor et al.(2024)]{BlankespoorDeHaanLi2024}
Blankespoor, E., deHaan, E., and Li, Q. (2024).
\newblock Generative AI in Financial Reporting.
\newblock \emph{SSRN Working Paper}, Abstract ID: 4986017.
\newblock Available at: \url{https://papers.ssrn.com/sol3/papers.cfm?abstract_id=4986017}.
\newblock \textit{Status: Working paper (last accessed March 2026). Not yet published in a journal to the author's knowledge.}

\bibitem[Callaway and Sant'Anna(2021)]{CallawayS2021}
Callaway, B. and Sant'Anna, P.H.C. (2021).
\newblock Difference-in-Differences with Multiple Time Periods.
\newblock \emph{Journal of Econometrics}, 225(2), 200--230.

\bibitem[Cao et al.(2023)]{CaoJiangYangZhang2023}
Cao, S., Jiang, W., Yang, B., and Zhang, A.L. (2023).
\newblock How to Talk When a Machine Is Listening: Corporate Disclosure in the Age of AI.
\newblock \emph{Review of Financial Studies}, 36(9), 3603--3642.
\newblock DOI: 10.1093/rfs/hhad021. Originally circulated as NBER Working Paper No.\ 27950 (2020).

\bibitem[Cao et al.(2024)]{CaoYangZhao2024}
Cao, Y., Liu, M., Qiu, J., and Zhao, R. (2024).
\newblock Information in Disclosing Emerging Technologies: Evidence from AI Disclosure.
\newblock \emph{SSRN Working Paper}, Abstract ID: 4987085.
\newblock Available at: \url{https://papers.ssrn.com/sol3/papers.cfm?abstract_id=4987085}.
\newblock \textit{Status: Working paper (last accessed March 2026).}

\bibitem[Chang and He(2025)]{ChangHe2025}
Chang, C.-Y. and He, X. (2025).
\newblock The Liabilities of Robots.Txt.
\newblock \emph{SSRN Working Paper}, Abstract ID: 5159436.
\newblock Available at: \url{https://papers.ssrn.com/sol3/papers.cfm?abstract_id=5159436}.
\newblock \textit{Status: Working paper (last accessed March 2026). Law review article, not finance/accounting journal.}

\bibitem[Li(2008)]{Li2008}
Li, F. (2008).
\newblock Annual Report Readability, Current Earnings, and Earnings Persistence.
\newblock \emph{Journal of Accounting and Economics}, 45(2--3), 221--247.
\newblock \textit{Note: This paper is sometimes cited as published in The Accounting Review. The correct journal is the Journal of Accounting and Economics. The Fog Index application to 10-K readability.}

\bibitem[Merton(1987)]{Merton1987}
Merton, R.C. (1987).
\newblock A Simple Model of Capital Market Equilibrium with Incomplete Information.
\newblock \emph{Journal of Finance}, 42(3), 483--510.

\bibitem[Reitinger(2025)]{Reitinger2025}
Reitinger, N. (2025).
\newblock Measured Failures of robots.txt: Legal and Empirical Insights for Regulating Robots.
\newblock \emph{SSRN Working Paper}, Abstract ID: 5555040.
\newblock Available at: \url{https://papers.ssrn.com/sol3/papers.cfm?abstract_id=5555040}.
\newblock \textit{Status: Working paper (last accessed March 2026). Empirical analysis of general websites, not focused on publicly traded firms.}

\end{thebibliography}

\end{document}
